\documentclass{amsart}
% For dealing with references we use the comment environment
\usepackage{verbatim}
\newenvironment{reference}{\comment}{\endcomment}
%\newenvironment{reference}{}{}
\newenvironment{slogan}{\comment}{\endcomment}
\newenvironment{history}{\comment}{\endcomment}

% For commutative diagrams we use Xy-pic
\usepackage[all]{xy}

% We use 2cell for 2-commutative diagrams.
\xyoption{2cell}
\UseAllTwocells

% We use multicol for the list of chapters between chapters
\usepackage{multicol}

% This is generally recommended for better output
\usepackage{lmodern}
\usepackage[T1]{fontenc}

% For cross-file-references

% Package for hypertext links:
\usepackage{hyperref}

% For any local file, say "hello.tex" you want to link to please

% Theorem environments.
%
\theoremstyle{plain}
\newtheorem{theorem}[subsection]{Theorem}
\newtheorem{proposition}[subsection]{Proposition}
\newtheorem{lemma}[subsection]{Lemma}

\theoremstyle{definition}
\newtheorem{definition}[subsection]{Definition}
\newtheorem{example}[subsection]{Example}
\newtheorem{exercise}[subsection]{Exercise}
\newtheorem{situation}[subsection]{Situation}

\theoremstyle{remark}
\newtheorem{remark}[subsection]{Remark}
\newtheorem{remarks}[subsection]{Remarks}

\numberwithin{equation}{subsection}

% Macros
%
\def\lim{\mathop{\mathrm{lim}}\nolimits}
\def\colim{\mathop{\mathrm{colim}}\nolimits}
\def\Spec{\mathop{\mathrm{Spec}}}
\def\Hom{\mathop{\mathrm{Hom}}\nolimits}
\def\Ext{\mathop{\mathrm{Ext}}\nolimits}
\def\SheafHom{\mathop{\mathcal{H}\!\mathit{om}}\nolimits}
\def\SheafExt{\mathop{\mathcal{E}\!\mathit{xt}}\nolimits}
\def\Sch{\mathit{Sch}}
\def\Mor{\mathop{\mathrm{Mor}}\nolimits}
\def\Ob{\mathop{\mathrm{Ob}}\nolimits}
\def\Sh{\mathop{\mathit{Sh}}\nolimits}
\def\NL{\mathop{N\!L}\nolimits}
\def\CH{\mathop{\mathrm{CH}}\nolimits}
\def\proetale{{pro\text{-}\acute{e}tale}}
\def\etale{{\acute{e}tale}}
\def\QCoh{\mathit{QCoh}}
\def\Ker{\mathop{\mathrm{Ker}}}
\def\Im{\mathop{\mathrm{Im}}}
\def\Coker{\mathop{\mathrm{Coker}}}
\def\Coim{\mathop{\mathrm{Coim}}}

% Boxtimes
%
\DeclareMathSymbol{\boxtimes}{\mathbin}{AMSa}{"02}

%
% Macros for moduli stacks/spaces
%
\def\QCohstack{\mathcal{QC}\!\mathit{oh}}
\def\Cohstack{\mathcal{C}\!\mathit{oh}}
\def\Spacesstack{\mathcal{S}\!\mathit{paces}}
\def\Quotfunctor{\mathrm{Quot}}
\def\Hilbfunctor{\mathrm{Hilb}}
\def\Curvesstack{\mathcal{C}\!\mathit{urves}}
\def\Polarizedstack{\mathcal{P}\!\mathit{olarized}}
\def\Complexesstack{\mathcal{C}\!\mathit{omplexes}}
% \Pic is the operator that assigns to X its picard group, usage \Pic(X)
% \Picardstack_{X/B} denotes the Picard stack of X over B
% \Picardfunctor_{X/B} denotes the Picard functor of X over B
\def\Pic{\mathop{\mathrm{Pic}}\nolimits}
\def\Picardstack{\mathcal{P}\!\mathit{ic}}
\def\Picardfunctor{\mathrm{Pic}}
\def\Deformationcategory{\mathcal{D}\!\mathit{ef}}

\setlength{\emergencystretch}{3em}
\begin{document}
\title[Stacks Project, Tag 00N1]{1310 --- Buchsbaum-Eisenbud rank and regular-sequence criterion for exact finite free complexes over Noetherian local rings}
\maketitle
\noindent Copyright (C) 2005--2025 Johan de Jong. Stacks Project authors. Source: \href{https://stacks.math.columbia.edu/tag/00N1}{Tag 00N1}.

\noindent Editorial difficulty note (not author text). Difficulty H2: The proof combines splitting of unit matrix entries, associated-prime rank detection and a nonzerodivisor induction with dimension-localization induction and the Peskine-Szpiro acyclicity mechanism. The two interacting induction arguments turn determinantal depth conditions into exactness, a substantial structural theorem beyond ordinary qualification algebra..

\noindent Editorial provenance note (not author text). History: excerpt from revision \texttt{a0444\allowbreak 6e57e\allowbreak c1fbc\allowbreak 252a8\allowbreak 71afc\allowbreak ec775\allowbreak 2fb28\allowbreak 07b14}, algebra.tex, lines 24833--24924. Reference presentation was adapted; original-author context and core support blocks were retained or added. The mathematical wording is unchanged. Licensed under GNU FDL 1.2 or later; no invariant sections or cover texts. See ../licenses/COPYING. Context blocks reproduce author definitions and supporting results; they are not additional counted problems.

\begin{situation}
\label{situation-complex}
Here $R$ is a ring, and we have a complex
$$
0
\to
R^{n_e}
\xrightarrow{\varphi_e}
R^{n_{e-1}}
\xrightarrow{\varphi_{e-1}}
\ldots
\xrightarrow{\varphi_{i + 1}}
R^{n_i}
\xrightarrow{\varphi_i}
R^{n_{i-1}}
\xrightarrow{\varphi_{i-1}}
\ldots
\xrightarrow{\varphi_1}
R^{n_0}
$$
In other words we require $\varphi_i \circ \varphi_{i + 1} = 0$
for $i = 1, \ldots, e - 1$.
\end{situation}

\begin{definition}
\label{definition-rank}
Let $R$ be a ring. Suppose that $\varphi : R^m \to R^n$ is a map
of finite free modules.
\begin{enumerate}
\item The {\it rank} of $\varphi$ is the maximal $r$ such that
$\wedge^r \varphi : \wedge^r R^m \to \wedge^r R^n$ is nonzero.
\item We let $I(\varphi) \subset R$ be the ideal generated by
the $r \times r$ minors of the matrix of $\varphi$, where $r$
is the rank as defined above.
\end{enumerate}
\end{definition}

\begin{lemma}
\label{lemma-add-trivial-complex}
Suppose $R$ is a ring. Let
$$
\ldots
\xrightarrow{\varphi_{i + 1}}
R^{n_i}
\xrightarrow{\varphi_i}
R^{n_{i-1}}
\xrightarrow{\varphi_{i-1}}
\ldots
$$
be a complex of finite free $R$-modules. Suppose that for some $i$
some matrix coefficient of the map $\varphi_i$ is invertible.
Then the displayed complex is isomorphic to the direct sum of a complex
$$
\ldots \to
R^{n_{i + 2}} \xrightarrow{\varphi_{i + 2}}
R^{n_{i + 1}} \to
R^{n_i - 1} \to
R^{n_{i - 1} - 1} \to
R^{n_{i - 2}} \xrightarrow{\varphi_{i - 2}}
R^{n_{i - 3}} \to
\ldots
$$
and the complex $\ldots \to 0 \to R \to R \to 0 \to \ldots$
where the map $R \to R$ is the identity map.
\end{lemma}

\begin{proof}
The assumption means, after a change of basis of
$R^{n_i}$ and $R^{n_{i-1}}$ that the first basis
vector of $R^{n_i}$ is mapped via $\varphi_i$ to the first basis
vector of $R^{n_{i-1}}$. Let $e_j$ denote the
$j$th basis vector of $R^{n_i}$ and $f_k$ the $k$th
basis vector of $R^{n_{i-1}}$. Write $\varphi_i(e_j)
= \sum a_{jk} f_k$. So $a_{1k} = 0$ unless $k = 1$
and $a_{11} = 1$. Change basis on $R^{n_i}$ again
by setting $e'_j = e_j - a_{j1} e_1$ for $j > 1$.
After this change of coordinates we have $a_{j1} = 0$
for $j > 1$. Note the image
of $R^{n_{i + 1}} \to R^{n_i}$ is contained in the
subspace spanned by $e_j$, $j > 1$. Note also
that $R^{n_{i-1}} \to R^{n_{i-2}}$ has to annihilate
$f_1$ since it is in the image. These conditions
and the shape of the matrix $(a_{jk})$ for $\varphi_i$
imply the lemma.
\end{proof}


\begin{lemma}
\label{lemma-exact-depth-zero-local}
In Situation \ref{situation-complex}. Suppose $R$ is
a local Noetherian ring with maximal ideal $\mathfrak m$.
Assume $\mathfrak m \in \text{Ass}(R)$, in other words
$R$ has depth $0$. Suppose that
$0 \to R^{n_e} \to R^{n_{e-1}} \to \ldots \to R^{n_0}$
is exact at $R^{n_e}, \ldots, R^{n_1}$.
Then the complex is isomorphic to a direct sum of trivial
complexes.
\end{lemma}

\begin{proof}
Pick $x \in R$, $x \not = 0$, with $\mathfrak m x = 0$.
Let $i$ be the biggest index such that $n_i > 0$.
If $i = 0$, then the statement is true. If
$i > 0$ denote $f_1$ the first basis vector of $R^{n_i}$.
Since $xf_1$ is not mapped to zero by
exactness of the complex we deduce that some matrix
coefficient of the map $R^{n_i} \to R^{n_{i - 1}}$
is not in $\mathfrak m$.
Lemma \ref{lemma-add-trivial-complex} then allows
us to decrease $n_e + \ldots + n_1$. Induction finishes the proof.
\end{proof}


\begin{lemma}
\label{lemma-trivial-case-exact}
In Situation \ref{situation-complex}, suppose the complex is
isomorphic to a direct sum of trivial complexes. Then
we have
\begin{enumerate}
\item the maps $\varphi_i$ have rank
$r_i = n_i - n_{i + 1} + \ldots + (-1)^{e-i-1} n_{e-1} + (-1)^{e-i} n_e$,
\item for all $i$, $1 \leq i \leq e - 1$ we have
$\text{rank}(\varphi_{i + 1}) + \text{rank}(\varphi_i) = n_i$,
\item each $I(\varphi_i) = R$.
\end{enumerate}
\end{lemma}

\begin{proof}
We may assume the complex is the direct sum of trivial
complexes. Then for each $i$ we can split the standard basis
elements of $R^{n_i}$ into those that map to a basis element
of $R^{n_{i-1}}$ and those that are mapped to zero (and these
are mapped onto by basis elements of $R^{n_{i + 1}}$ if $i > 0$).
Using descending
induction starting with $i = e$ it is easy to prove that there
are $r_{i + 1}$-basis elements of $R^{n_i}$ which are mapped
to zero and $r_i$ which are mapped to basis elements of
$R^{n_{i-1}}$. From this the result follows.
\end{proof}


\begin{lemma}
\label{lemma-div-x-exact-one-less}
In Situation \ref{situation-complex}. Suppose $R$ is
a local ring with maximal ideal $\mathfrak m$.
Suppose that $0 \to R^{n_e} \to R^{n_{e-1}} \to \ldots \to R^{n_0}$
is exact at $R^{n_e}, \ldots, R^{n_1}$.
Let $x \in \mathfrak m$ be a nonzerodivisor. The complex
$0 \to (R/xR)^{n_e} \to \ldots \to (R/xR)^{n_1}$
is exact at $(R/xR)^{n_e}, \ldots, (R/xR)^{n_2}$.
\end{lemma}

\begin{proof}
Denote $F_\bullet$ the complex with terms $F_i = R^{n_i}$
and differential given by $\varphi_i$. Then we have a short
exact sequence of complexes
$$
0 \to F_\bullet \xrightarrow{x} F_\bullet \to F_\bullet/xF_\bullet \to 0
$$
Applying the snake lemma we get a long exact sequence
$$
H_i(F_\bullet) \xrightarrow{x} H_i(F_\bullet) \to
H_i(F_\bullet/xF_\bullet) \to H_{i - 1}(F_\bullet)
\xrightarrow{x} H_{i - 1}(F_\bullet)
$$
The lemma follows.
\end{proof}


\begin{lemma}[Acyclicity lemma]
\label{lemma-acyclic}
\begin{reference}
\href{https://stacks.math.columbia.edu/bibliography}{Bibliography: Peskine-Szpiro (Lemma 1.8)}
\end{reference}
Let $R$ be a local Noetherian ring.
Let $0 \to M_e \to M_{e-1} \to \ldots \to M_0$
be a complex of finite $R$-modules.
Assume $\text{depth}(M_i) \geq i$.
Let $i$ be the largest index such that the complex is
not exact at $M_i$. If $i > 0$ then
$\Ker(M_i \to M_{i-1})/\Im(M_{i + 1} \to M_i)$
has depth $\geq 1$.
\end{lemma}

\begin{proof}
Let $H = \Ker(M_i \to M_{i-1})/\Im(M_{i + 1} \to M_i)$ be the
cohomology group in question.
We may break the complex into short exact sequences
$0 \to M_e \to M_{e-1} \to K_{e-2} \to 0$,
$0 \to K_j \to M_j \to K_{j-1} \to 0$, for $i + 2 \leq j \leq e-2 $,
$0 \to K_{i + 1} \to M_{i + 1} \to B_i \to 0$,
$0 \to K_i \to M_i \to M_{i-1}$, and
$0 \to B_i \to K_i \to H \to 0$.
We proceed up through these complexes to
prove the statements about depths, repeatedly using
Lemma \href{https://stacks.math.columbia.edu/tag/00LX}{lemma depth in ses (Tag 00LX)}.
First of all, since $\text{depth}(M_e) \geq e$,
and $\text{depth}(M_{e-1}) \geq e-1$ we deduce
that $\text{depth}(K_{e-2}) \geq e - 1$. At this point the
sequences $0 \to K_j \to M_j \to K_{j-1} \to 0$ for $i + 2 \leq j \leq e-2 $
imply similarly that $\text{depth}(K_{j-1}) \geq j$ for
$i + 2 \leq j \leq e-2$. The sequence
$0 \to K_{i + 1} \to M_{i + 1} \to B_i \to 0$
then shows that $\text{depth}(B_i) \geq i + 1$. The sequence
$0 \to K_i \to M_i \to M_{i-1}$ shows that $\text{depth}(K_i) \geq 1$
since $M_i$ has depth $\geq i \geq 1$ by assumption.
The sequence $0 \to B_i \to K_i \to H \to 0$ then
implies the result.
\end{proof}


\begin{proposition}
\label{proposition-what-exact}
\begin{reference}
\href{https://stacks.math.columbia.edu/bibliography}{Bibliography: WhatExact (Corollary 1)}
\end{reference}
In Situation \ref{situation-complex}, suppose $R$ is
a local Noetherian ring. The following are equivalent
\begin{enumerate}
\item $0 \to R^{n_e} \to R^{n_{e-1}} \to \ldots \to R^{n_0}$
is exact at $R^{n_e}, \ldots, R^{n_1}$, and
\item for all $i$, $1 \leq i \leq e$
the following two conditions are satisfied:
\begin{enumerate}
\item $\text{rank}(\varphi_i) = r_i$ where
$r_i = n_i - n_{i + 1} + \ldots + (-1)^{e-i-1} n_{e-1} + (-1)^{e-i} n_e$,
\item $I(\varphi_i) = R$, or $I(\varphi_i)$ contains a
regular sequence of length $i$.
\end{enumerate}
\end{enumerate}
\end{proposition}

\begin{proof}
If for some $i$ some matrix coefficient of $\varphi_i$
is not in $\mathfrak m$, then we apply Lemma \ref{lemma-add-trivial-complex}.
It is easy to see that the proposition for a complex and
for the same complex with a trivial complex added to it
are equivalent. Thus we may assume that all matrix entries
of each $\varphi_i$ are elements of the maximal ideal.
We may also assume that $e \geq 1$.

\medskip\noindent
Assume the complex is exact at $R^{n_e}, \ldots, R^{n_1}$.
Let $\mathfrak q \in \text{Ass}(R)$.
Note that the ring $R_{\mathfrak q}$ has depth $0$
and that the complex remains exact after localization at $\mathfrak q$.
We apply Lemmas \ref{lemma-exact-depth-zero-local} and
\ref{lemma-trivial-case-exact} to the localized complex
over $R_{\mathfrak q}$. We conclude that
$\varphi_{i, \mathfrak q}$ has rank $r_i$ for all $i$.
Since $R \to \bigoplus_{\mathfrak q \in \text{Ass}(R)} R_\mathfrak q$
is injective (Lemma \href{https://stacks.math.columbia.edu/tag/0311}{lemma zero at ass zero (Tag 0311)}), we conclude that
$\varphi_i$ has rank $r_i$ over $R$ by the definition of rank as given
in Definition \ref{definition-rank}. Therefore we see that
$I(\varphi_i)_\mathfrak q = I(\varphi_{i, \mathfrak q})$
as the ranks do not change. Since all of the ideals
$I(\varphi_i)_{\mathfrak q}$, $e \geq i \geq 1$
are equal to $R_{\mathfrak q}$ (by the lemmas referenced above)
we conclude none of the ideals $I(\varphi_i)$ is contained in $\mathfrak q$.
This implies that $I(\varphi_e)I(\varphi_{e-1})\ldots I(\varphi_1)$
is not contained in any of the associated primes
of $R$. By Lemma \href{https://stacks.math.columbia.edu/tag/00DS}{lemma silly (Tag 00DS)} we may choose
$x \in I(\varphi_e)I(\varphi_{e - 1})\ldots I(\varphi_1)$,
$x \not \in \mathfrak q$ for all $\mathfrak q \in \text{Ass}(R)$.
Observe that $x$ is a nonzerodivisor (Lemma \href{https://stacks.math.columbia.edu/tag/00LD}{lemma ass zero divisors (Tag 00LD)}).
According to Lemma \ref{lemma-div-x-exact-one-less}
the complex $0 \to (R/xR)^{n_e} \to \ldots \to (R/xR)^{n_1}$ is exact
at $(R/xR)^{n_e}, \ldots, (R/xR)^{n_2}$. By induction
on $e$ all the ideals $I(\varphi_i)/xR$ have a regular
sequence of length $i - 1$. This proves that $I(\varphi_i)$
contains a regular sequence of length $i$.

\medskip\noindent
Assume (2)(a) and (2)(b) hold. We will prove that (1) holds by
induction on $\dim(R)$. If $\dim(R) = 0$, then we must have
$I(\varphi_i) = R$ for $1 \leq i \leq e$ by (2)(b). Since the
coefficients of $\varphi_i$ are contained in the maximal ideal
this can happen only if $r_i = 0$ for all $i$. By (2)(a) we
conclude that $e = 0$ and (1) holds. Assume $\dim(R) > 0$.
We claim that for any prime $\mathfrak p \subset R$ conditions
(2)(a) and (2)(b) hold for the complex
$0 \to R_\mathfrak p^{n_e} \to R_\mathfrak p^{n_{e - 1}} \to \ldots \to
R_\mathfrak p^{n_0}$ with maps $\varphi_{i, \mathfrak p}$
over $R_\mathfrak p$. Namely, since $I(\varphi_i)$ contains a
nonzero divisor, the image of $I(\varphi_i)$ in $R_\mathfrak p$
is nonzero. This implies that the rank of $\varphi_{i, \mathfrak p}$
is the same as the rank of $\varphi_i$: the rank as defined above
of a matrix $\varphi$ over a ring $R$ can only drop when passing
to an $R$-algebra $R'$ and this happens if and only if $I(\varphi)$
maps to zero in $R'$. Thus (2)(a) holds. Having said this
we know that $I(\varphi_{i, \mathfrak p}) = I(\varphi_i)_\mathfrak p$
and we see that (2)(b) is preserved under localization as well.
By induction on the dimension of $R$ we may assume the complex
is exact when localized at any nonmaximal prime $\mathfrak p$ of $R$.
Thus $\Ker(\varphi_i)/\Im(\varphi_{i + 1})$ has support contained in
$\{\mathfrak m\}$ and hence if nonzero has depth $0$.
As $I(\varphi_i) \subset \mathfrak m$ for all $i$ because
of what was said in the first paragraph of the proof, we
see that (2)(b) implies $\text{depth}(R) \geq e$.
By Lemma \ref{lemma-acyclic} we see
that the complex is exact at $R^{n_e}, \ldots, R^{n_1}$
concluding the proof.
\end{proof}
\end{document}
